\section{Project and Context}
\label{chap:project_context}

This master’s project is the continuation of a specialization project conducted in the fall of 2020 (Mykland and Tyldum, 2020). The written assignment for both the specialization project and master’s thesis was as follows:

\blockquote{
    \textbf{[Game technology] Alternative games} \\
    \textit{In this project, the goal is to prototype an innovative game and test this game on users. The innovation can be in the type of gameplay the game provides, how it combines various game genres, what technology is used to control or play the game, how the social interaction between players is supported, the purpose of the game, etc.}
}

Due to this assignment’s large scope, the project was separated into four phases. Two of which were conducted during the specialization project and two in this master’s thesis.

The first phase, conducted in the specialization project, included a study of game theory, alternative games, game technologies, and existing concepts. This study helped us develop a model for creating new alternative game concepts. The second phase was to explore and develop new game concepts that generate value for individual players and other society members. This phase utilized the model developed in phase one and led to the development of the game concept Hover. In this report, the relevant chapters from the specialization project are included in Part II. 

In this master’s thesis, we proceed with the last two phases of the project. The third phase  was to implement a prototype of the game concept, and the fourth phase was to study and evaluate the benefits of said concept through user testing. The game design and rough implementation plan created in the specialization project enabled us to start phase game development in part three quickly. The prototype took three months to develop, with a total of 36500 lines of code.

The planning of phase four was conducted in parallel with phase three. It involved defining which game metrics to test, how and when they should be tested, who to test it on, how to measure and collect data, and planning the game’s distribution. The research scheme was created in line with guidelines provided by NTNU and the Norwegian Centre for Research Data (NSD), who also approved the project.

The user testing in phase four lasted two weeks, and participants were asked to answer a survey afterward. The results were then processed and analyzed to conclude phase four.

The work presented in this master’s thesis is our contribution to the field of alternative games.
